\chapter{Background and Related Work}
\label{ch:related-work}

This chapter reviews the research that supports the three perception
sub-tasks studied in this thesis, together with the work on integrated
analysis systems that motivates their composition. The review is organised
thematically rather than paper by paper. Section~\ref{sec:rw-detection}
establishes the two dominant models for localising objects in images,
bounding-box regression and heatmap regression, and explains why the latter
is better suited to small, fast targets. Section~\ref{sec:rw-ball} surveys
ball tracking in sports video, with particular attention to the TrackNet
family that this thesis benchmarks. Section~\ref{sec:rw-court} covers court
and field registration, spanning keypoint-detection backbones and the
geometry of homography estimation. Section~\ref{sec:rw-event} reviews event
detection from trajectories, including temporal convolutional networks and
gradient-boosted decision trees. Section~\ref{sec:rw-systems} situates the
work within integrated tennis systems and automated officiating, and
Section~\ref{sec:rw-gap} synthesises the preceding material into the research
gap that this thesis addresses.

\section{Object Detection and Small-Object Localisation}
\label{sec:rw-detection}

\subsection{Bounding-box detection}
\label{sec:rw-bbox}

Modern object detection was reshaped by the region-based convolutional
neural network (R-CNN), which combined bottom-up region proposals with
convolutional features and demonstrated a large accuracy gain over methods
built on hand-crafted features \cite{girshick2014rcnn}. The original design
was computationally expensive because it processed every proposal
independently; Fast R-CNN removed this redundancy by sharing one
convolutional pass over the whole image and pooling features per region
\cite{girshick2015fastrcnn}, and Faster R-CNN replaced the external
proposal step with a learned region proposal network, yielding a near
end-to-end two-stage detector \cite{ren2017fasterrcnn}. In parallel,
one-stage detectors removed the proposal stage altogether: the single shot
multibox detector (SSD) predicted boxes directly from multiple feature maps
\cite{liu2016ssd}, while the You Only Look Once (YOLO) family framed
detection as a single regression over a spatial grid
\cite{redmon2016yolo}. This last design accepted a modest accuracy
reduction in return for substantially higher throughput. Successive YOLO revisions refined the backbone and the
multi-scale prediction head \cite{redmon2018yolov3}, and the line has since
continued through widely used engineering releases that, at the time of
writing, are documented as software rather than peer-reviewed publications
\cite{jocher2024yolo11}. Two architectural ideas from this lineage are
especially relevant to small objects. Feature pyramid networks (FPNs) build a
top-down feature hierarchy with lateral connections so that small objects
remain represented at high spatial resolution \cite{lin2017fpn}, and the
focal loss addresses the extreme foreground-background imbalance that
otherwise overwhelms dense one-stage training \cite{lin2017focalloss}.

Despite these advances, detecting small objects remains difficult. A recent
survey identifies the loss of spatial detail through repeated downsampling,
the scarcity of discriminative appearance cues, and the imbalance between
tiny foreground regions and large backgrounds as persistent obstacles
\cite{cheng2023smallobjectsurvey}. A tennis ball is an acute instance of this
problem. It occupies only a few pixels and is frequently blurred by motion,
and because its bounding box is barely larger than a point, the
box-regression objective conveys little information beyond a location.

\subsection{Heatmap regression}
\label{sec:rw-heatmap}

An alternative to box regression represents a target as a spatial probability
map, in which a Gaussian peak marks the object centre and the predicted
coordinate is recovered by locating the maximum response. This formulation
originated in human pose estimation, where convolutional pose machines
refined per-joint belief maps through successive stages \cite{wei2016cpm} and
stacked hourglass networks repeatedly pooled and upsampled features to
aggregate evidence across scales \cite{newell2016hourglass}. The idea was
later generalised to object detection by representing each object as a single
centre point in a heatmap, which removed the need for anchor boxes and
non-maximum suppression \cite{zhou2019centernet}. Heatmap regression is
attractive for small, fast targets for several reasons. It supplies dense
per-pixel supervision, and it localises a target sub-pixel through the shape
of the response rather than through a discrete box. It also accommodates the
absence of a target naturally, as a map with no significant peak. These
properties make it the dominant formulation for both the ball-tracking and
court-keypoint sub-tasks studied in this thesis, and they motivate
comparing heatmap-based trackers with a single bounding-box detector when
those sub-tasks are taken up in later chapters.

\section{Ball Tracking in Sports Video}
\label{sec:rw-ball}

\subsection{Tracking by detection}
\label{sec:rw-tbd}

A common strategy for following an object through a video is tracking by
detection, in which a per-frame detector is coupled with a temporal data
association step. Simple online and realtime tracking (SORT) associates
detections across frames using a Kalman-filter motion model and the
Hungarian assignment algorithm, achieving high speed with bounding-box
overlap as the sole association cue \cite{bewley2016sort}, and DeepSORT
augments this with a learned appearance descriptor to reduce identity switches
through occlusion \cite{wojke2017deepsort}. These methods are effective for
pedestrians and vehicles, but they rest on assumptions that a tennis ball
violates: the ball has almost no distinctive appearance to embed, it moves
fast enough to blur or skip across the frame, and it is regularly occluded by
players, the net, and court lines. Single-frame detection followed by
geometric association is therefore fragile in this setting, which has
motivated trackers that reason over several frames jointly.

\subsection{The TrackNet family}
\label{sec:rw-tracknet}

TrackNet introduced a heatmap-based, multi-frame approach designed
specifically for small and fast sports objects \cite{huang2019tracknet}.
Its encoder-decoder network ingests three consecutive frames, stacked into a
nine-channel input, and outputs a Gaussian heatmap locating the ball, so that
motion across the input frames provides the temporal cue absent from a single
image. This original design follows a multiple-in single-out pattern, in which
several input frames yield one output heatmap. TrackNetV2 reorganised the
network into a multiple-in multiple-out architecture that predicts a heatmap
for every input frame in one pass, reducing redundant computation and
improving throughput while retaining accuracy on shuttlecock data
\cite{sun2020tracknetv2}. TrackNetV3 added two refinement components on top of
this backbone: an auxiliary background-aware module to improve robustness under
visual interference, and an inpainting-based trajectory-rectification module
that repairs missed detections by reasoning over the predicted path
\cite{chen2023tracknetv3}. TrackNetV4 introduced a lightweight motion-attention
mechanism, in which inter-frame difference maps are fused with the backbone
features so that the network explicitly attends to moving-ball regions, with
reported gains under occlusion and low visibility \cite{raj2025tracknetv4}.
The lineage continues to be extended through further variants proposed as
preprints \cite{tang2025tracknetv5}, but this thesis treats the published
members of the family as its primary reference points.

Beyond the TrackNet line, related work has tackled ball and event analysis in
adjacent racket and table sports. TTNet performs joint ball detection, semantic
segmentation, and temporal event spotting for table tennis within a single
multi-task network \cite{voeikov2020ttnet}. This shows that detection and
event recognition can share a representation. Patch-wise convolutional
classifiers have also been applied to tennis ball detection in broadcast video,
linking local detections into trajectories \cite{reno2018cnntennis}. Taken
together, this literature establishes multi-frame heatmap regression as the
prevailing paradigm for localising a small, fast ball, and frames the
comparison undertaken in Chapter~\ref{ch:ball} between successive members of
the TrackNet family and a representative single-frame detector.

\section{Court Keypoint Detection and Geometric Registration}
\label{sec:rw-court}

\subsection{Keypoint detection backbones}
\label{sec:rw-keypoint}

Recovering the geometry of the playing surface begins with localising a fixed
set of court landmarks, a task closely related to human-pose keypoint
estimation. The heatmap-regression backbones reviewed in
Section~\ref{sec:rw-heatmap} transfer directly to this problem. A widely used
baseline appends a few deconvolutional layers to a deep residual network
(ResNet) backbone \cite{he2016resnet} to produce keypoint heatmaps from
high-level features, a design valued for its simplicity \cite{xiao2018simplebaselines}.
In contrast, the high-resolution network (HRNet) maintains high-resolution
representations throughout the network and fuses parallel multi-resolution
streams, which improves the spatial precision of the predicted keypoints
\cite{sun2019hrnet}. Because court detection must often run within a larger
real-time pipeline, efficient backbones are also of interest: the MobileNet
family uses inverted-residual blocks with linear bottlenecks \cite{sandler2018mobilenetv2}
and architecture search with squeeze-and-excitation refinements
\cite{howard2019mobilenetv3} to deliver competitive accuracy at a small
fraction of the parameter count. Whether such lightweight backbones suffice for
accurate court-keypoint localisation is one of the open questions this thesis
sets out to answer.

\subsection{Sports-field registration and homography}
\label{sec:rw-homography}

Localised keypoints are useful for analysis only once they are related to a
metric model of the court. For a planar surface viewed by a single camera,
this relationship is a homography, a projective transformation between the
image plane and the court plane whose theory is set out in the standard
reference on multiple-view geometry \cite{hartley2004mvg}. Estimating the
homography from automatically detected keypoints requires robustness to
mislocalised points, which is classically provided by random sample consensus
(RANSAC), an iterative scheme that fits a model to randomly sampled minimal
subsets and retains the hypothesis with the largest consensus set
\cite{fischler1981ransac}. Learned variants such as differentiable RANSAC make
the hypothesis-selection step trainable end to end \cite{brachmann2017dsac}.
Within sports specifically, field registration has been framed as inference in
a deep structured model that combines semantic cues with geometric priors
\cite{homayounfar2017sportsfield}, as iterative optimisation through a learned
error function \cite{jiang2020sportsfield}, and as dense keypoint detection on
a regular grid coupled with feature-warp refinement to remain robust where
distinctive corners are absent \cite{nie2021sportsfield}. The court pipeline in
Chapter~\ref{ch:court} follows the detect-then-fit pattern common to this
literature: predict court keypoints with a heatmap backbone, then fit a
homography with a robust estimator to map the image into court-metric
coordinates.

\section{Event Detection from Trajectories}
\label{sec:rw-event}

\subsection{Temporal modelling of sequences}
\label{sec:rw-temporal}

Identifying the discrete frames at which a bounce occurs is a temporal
pattern-recognition problem over a one-dimensional signal derived from the
ball trajectory. Although recurrent networks were long the default for
sequence modelling, convolutional alternatives have proved competitive.
Dilated causal convolutions, popularised for raw-audio generation by WaveNet,
expand the receptive field exponentially with depth without increasing the
parameter count \cite{vandenoord2016wavenet}. Long temporal contexts can
therefore be captured efficiently. Temporal convolutional networks (TCNs) build on
this principle for discriminative tasks: they have been applied to fine-grained
action segmentation and detection in video \cite{lea2017tcn}, and a systematic
comparison found that generic TCNs match or exceed recurrent networks across a
range of sequence-modelling benchmarks while being simpler to train
\cite{bai2018tcn}. These properties make a dilated TCN a natural candidate for
classifying windows of trajectory features into bounce events.

\subsection{Gradient-boosted decision trees}
\label{sec:rw-gbm}

A problem is naturally expressed through engineered per-frame features,
gradient-boosted decision trees offer a strong and data-efficient alternative
to deep models. Gradient boosting fits an additive ensemble of trees by
successively approximating the negative gradient of a loss function
\cite{friedman2001gbm}. Modern implementations differ mainly in how they
control overfitting and scale to large data: XGBoost introduced a regularised
objective and a scalable system design \cite{chen2016xgboost}, LightGBM
accelerated training through histogram-based splits and feature bundling
\cite{ke2017lightgbm}, and CatBoost reduced the prediction bias of ordinary
boosting through ordered boosting \cite{prokhorenkova2018catboost}.
Histogram-based gradient boosting is also available as a standard component of
general machine-learning toolkits \cite{pedregosa2011sklearn}. Such models
operate directly on the kinematic and shape features computed from a trajectory
and require comparatively little tuning, which makes a gradient-boosted
classifier a natural baseline among the bounce detectors this thesis compares.

\subsection{Bounce and hit detection in tennis}
\label{sec:rw-bounce}

Compared with ball localisation, the detection of bounce and hit events from a
tennis trajectory has received comparatively little dedicated attention, and
existing approaches range from hand-built kinematic heuristics that look for
characteristic reversals in vertical motion to learned classifiers over
trajectory features. Multi-task systems such as TTNet show that event spotting
can be learned jointly with detection \cite{voeikov2020ttnet}, but such work
generally targets a different sport and a different event vocabulary. The
relative scarcity of controlled comparisons between tree-based and
temporal-convolutional detectors for tennis bounces, evaluated on a common
trajectory under a common tolerance, is one of the gaps that
Chapter~\ref{ch:bounce} addresses.

\section{Integrated Tennis Systems and Automated Officiating}
\label{sec:rw-systems}

The most prominent application of tennis computer vision is automated
officiating. The Hawk-Eye system reconstructs the three-dimensional ball
trajectory from several synchronised, calibrated cameras and estimates the
bounce position to adjudicate line calls \cite{owens2003hawkeye}. Its
commercial success demonstrates the value of vision-based measurement, yet
scholarly analysis has cautioned that the confident, deterministic
visualisations such systems present can mask the measurement uncertainty
inherent in any reconstruction of ball position \cite{collins2008hawkeye}.
Equally important for the present work, these systems depend on calibrated
multi-camera installations that are unavailable outside elite venues, which
motivates accessible, vision-only analysis of ordinary monocular video.

More broadly, surveys of computer vision for sport document a wide range of
perception and analysis tasks, from player and ball tracking to tactical
understanding, and note that components are typically developed and reported in
isolation \cite{thomas2017cvsports, zhao2023sportssurvey}. The literature
contains strong individual models for ball tracking, court registration, and
event detection, but comparatively few studies compose these components into a
single monocular pipeline whose end-to-end behaviour and deployability can be
examined, which is precisely the integration studied in
Chapter~\ref{ch:integrated}.

\section{Research Gap and Positioning}
\label{sec:rw-gap}

Three observations emerge from the preceding review. First, within each of the
three perception sub-tasks, a wide range of architectures has been proposed,
but they are commonly reported on heterogeneous datasets and private
train/test splits, and most studies evaluate a single proposed model rather
than placing competing designs on a common footing. This makes it difficult to
draw like-for-like conclusions about the accuracy, speed, and robustness
trade-offs between architectures within a sub-task. Second, the heatmap
formulation that recurs across ball tracking and court detection, and the
temporal and tree-based models that recur across event detection, are seldom
compared head to head under a fixed protocol on the same data. Third, and to
the best of the author's knowledge, comparatively few works compose the
resulting components into a single end-to-end monocular pipeline whose
practical deployability, including model export and per-frame latency, is
assessed, and fewer still close the loop to a geometric line-call decision from
ordinary single-camera video.

This thesis is positioned against these gaps. It conducts a controlled,
within-task comparison for each of the three sub-tasks under a deterministic
evaluation protocol, and it composes the strongest model from each into one
reproducible pipeline that maps detected bounces through an estimated
homography to a qualitative in/out line call. The protocol, datasets, and
shared metric vocabulary that make these comparisons reproducible are described
next, in Chapter~\ref{ch:methodology}.
